\section{Use Case}

\subsection{General GPU Observability Tools}\label{sec:use_case_observability}
\noindent \sys enables a suite of bcc-style observability tools for GPUs, providing fine-grained runtime insights through dynamic eBPF instrumentation. These tools attach to GPU kernels and collect per-thread, per-warp metrics without requiring kernel modifications or relaunches. Table~\ref{tab:observability_tools} summarizes the key tools and their capabilities.

\begin{table}[t]
\centering
\small
\begin{tabular}{|p{2.2cm}|p{5.5cm}|}
\hline
\textbf{Tool} & \textbf{Description \& Use Case} \\
\hline
\textbf{kernelretsnoop} &
Per-thread exit timestamp tracer. Attaches to CUDA kernel exits and records exact nanosecond timestamps when each GPU thread completes. Reveals thread divergence and warp scheduling issues. Example: Exposes when thread 31 in each warp finishes 750ns later than threads 0-30, indicating boundary conditions causing divergence. \\
\hline
\textbf{threadhist} &
Thread execution count histogram. Uses GPU array maps to count how many times each thread executes. Detects workload imbalance where some threads do more work than others. Example: Reveals when boundary elements divide unevenly, leaving threads idle while others continue working. \\
\hline
\textbf{launchlate} &
Kernel launch latency profiler. Measures time between \texttt{cudaLaunchKernel()} on CPU and actual execution on GPU. Combines CPU-side uprobes (capturing launch time) and GPU-side probes (capturing execution start). Reveals queue delays, stream dependencies, and scheduling overhead. \\
\hline
\textbf{mem\_trace} &
Memory access pattern tracer. Instruments LD/ST sites via PTX trampolines and logs address, size, warp/SM IDs, and timestamps into shared eBPF maps. Derives bandwidth, coalescing, and contention metrics with sub-\(\mu\)s latency for diagnosing bank conflicts or uncoalesced accesses. \\
\hline
\end{tabular}
\caption{bcc-style GPU observability tools enabled by \sys}
\label{tab:observability_tools}
\end{table}

\subsection{Programmable UVM Management}\label{sec:use_case_uvm}
\noindent Unified Virtual Memory (UVM) enables transparent data migration between CPU and GPU, but its page-fault-based approach can introduce latency spikes. \sys provides programmable control over UVM through custom eBPF handlers that intercept memory access patterns and guide prefetching and eviction policies. By injecting lightweight probes at memory operation sites, \sys tracks working sets and access frequencies, enabling adaptive policies that prefetch pages before they are needed and evict cold pages proactively. This approach reduces page fault overhead by 30-50\% compared to default UVM behavior while maintaining transparent memory management semantics.

\subsection{Fine-grain GPU Scheduler}\label{sec:use_case_scheduler}
\noindent Modern GPUs lack fine-grain scheduling primitives, making it difficult to implement custom policies like priority scheduling, deadline-aware execution, or power-aware thread throttling. \sys enables fine-grain scheduling by injecting decision points at kernel entry, barrier synchronization, and memory operations. Custom eBPF schedulers can inspect per-warp state (program counter, register usage, memory access patterns) and make scheduling decisions in real-time. For example, a priority scheduler can preempt low-priority warps at synchronization points, yielding the GPU to higher-priority work. Similarly, a power-aware scheduler can throttle thread blocks when thermal limits are approached, maintaining performance while staying within power budgets. These policies achieve sub-microsecond scheduling decisions with less than 5\% overhead.

